\section{Variance after the inverse estimate}\label{sec:variance}

The inverse estimate leaves a fixed finite list of rational alignments in each
candidate box. Fix one rational and every endpoint factor except the
isolated prime. We average over the whole isolated bin by temporarily
imposing the aligned hit class at each prime; at the primes that actually
align, this gives the original hit progression. We show that a positive
fraction of these progressions can have large small-prime sieve errors
only when two arithmetic parameters of the family are small. One controls
the nonaligned factors; the other controls the translation of the
progression after the aligned factors are removed.
The auxiliary probability models used below are explicit dominating measures;
we do not assume that the residue classes selected by the actual primes are
independent.

\subsection{Arithmetic coordinates}

Fix a candidate box from \cref{prop:regular-boxes}, and write
\(\mathcal P\) for its isolated prime bin, \(R\) for the lower endpoint of
that bin, and \(n_p=\#\mathcal P\). Thus
\[
 z^{\alpha_{\min}}\le R\le z^{\alpha_{\max}},
 \qquad n_p\asymp \frac{\xi R}{\log R}.
\]
The constants in comparisons of this form can be taken independent of
sufficiently small fixed \(\xi\); the threshold on \(z\) may depend on \(\xi\).
Fix all prime choices in a final compact endpoint except the isolated
\(p\in\mathcal P\), and write
\[
 d=pq_f.
\]
Fix also one reduced rational \(A/D\), \(D>0\), from the candidate list in
\cref{prop:inverse}. All of \(q_f,A,D\) remain fixed as \(p\) varies.
As before, let \(q_{f,E}\) be the product of the nonaligned prime factors
of \(q_f\), and put \(q_{f,A}=q_f/q_{f,E}\).

If the structured-parent conclusion of \cref{prop:inverse} holds, the
product of the candidate edge and the later factors is at most a fixed
power of \(w\), by the candidate parent-gap bound. Hence, for large \(z\),
\begin{equation}\label{eq:variance-structure-size}
 H:=Dq_{f,E}\le R\mathcal Z^{11},
 \qquad
 |A|\le S_qR\mathcal Z^{10}.
\end{equation}
Here \(q\) is the product of the other factors at the candidate parent,
\(S_q\) is the corresponding product of bin upper endpoints, and
\[
 q_f\ge q,\qquad S_q/q\le\mathcal Z^{o(1)}.
\]
The additional factor \(\mathcal Z\) in the bound for \(H\) absorbs any
fixed power of \(w\).

Let \(b\in[1,q_f]\) be the integer CRT representative of the hit
conditions at the primes of \(q_f\), and define
\begin{equation}\label{eq:variance-effective}
 C_0=\frac{Db-A}{q_{f,A}},
 \qquad T=\max\left(H,\frac{|C_0|}{R}\right).
\end{equation}
The numerator is divisible by every aligned factor of \(q_f\), so \(C_0\)
is an integer. To see the role of these parameters, write the hit
progression at \(q_f\) as \(n=b+q_f v\). Then
\[
 Dn-A=q_{f,A}(C_0+Hv).
\]
Thus removing the aligned factors leaves an arithmetic progression of
step \(H\) and initial value \(C_0\). Alignment at the isolated prime
will require this remaining value to be divisible by \(p\).
The scale \(T\) measures both its step and the size of its initial value
after division by a prime in the bin \(p\asymp R\).

We also have
\begin{equation}\label{eq:variance-unit}
 (C_0,H)=1.
\end{equation}
Indeed, \(q_{f,A}\) is coprime to \(Dq_{f,E}\). At a prime \(t\mid D\),
reducedness gives \(Db-A\equiv-A\not\equiv0\pmod t\); at a prime
\(t\mid q_{f,E}\) not dividing \(D\), nonalignment gives
\(Db-A\equiv Da_t-A\not\equiv0\pmod t\).
This also covers common prime divisors of \(D\) and \(q_{f,E}\).
Prime powers in \(D\) cause no further restriction.
In particular, \(C_0=0\) can occur only when \(H=1\).
The height bound gives
\[
 \frac{|C_0|}{RH}
 =\frac{|Db-A|}{RDq_f}
 \le \frac1R+\frac{|A|}{RDq_f}
 \le \frac1R+\frac{S_q}{q}\mathcal Z^{10},
\]
and therefore
\begin{equation}\label{eq:variance-effective-ratio}
 T/H\le\mathcal Z^{12}
\end{equation}
for sufficiently large \(z\).

For any \(p\in\mathcal P\) not dividing \(H\), temporarily impose the
aligned class at \(p\), whether or not its actual prescribed class aligns.
Since \(p\nmid q_f\), the additional hit condition is
\(C_0+Hv\equiv0\pmod p\). There is a unique
\(k\in\{0,\ldots,p-1\}\) satisfying this congruence, and all its
solutions have the form \(v=k+pj\). Define
\[
 m_0=\frac{C_0+Hk}{p}.
\]
Equivalently, \(m_0\) is the unique integer satisfying
\begin{equation}\label{eq:variance-m-coordinate}
 m_0\in[C_0/p,C_0/p+H),
 \qquad pm_0\equiv C_0\pmod H.
\end{equation}
The interval has length \(H\), and the congruence fixes one residue
modulo \(H\), which also verifies uniqueness in these coordinates.
The resulting simultaneous hit progression is exactly
\begin{equation}\label{eq:variance-progression}
 n=b+q_f(k+pj),\qquad j\ge0,
 \qquad
 Dn-A=pq_{f,A}(m_0+Hj).
\end{equation}
Its initial representative belongs to \([1,pq_f]\), so the allowable
\(j\)'s in \(1\le n\le Y\) form an initial integer interval of length
\(Y/(pq_f)+O(1)\).

An included endpoint has \(J_d\ge w^{a_\star}\). The within-box movement
of \(\log_w J_d\) is at most \(\Delta\), by
\cref{prop:regular-boxes}. Taking \(\Delta<a_\star/4\), every endpoint
choice in the full box therefore has \(J_d\ge w^{3a_\star/4}\).
Set
\[
 J=\left\lfloor\frac{Y}{(1+\xi)Rq_f}\right\rfloor,
 \qquad \mathcal J=\{0,\ldots,J-1\}.
\]
Since \(p\le(1+\xi)R\) and the initial representative lies in
\([1,pq_f]\), this interval is contained in every allowed
\(j\)-interval. Its cardinality satisfies
\begin{equation}\label{eq:variance-common-interval}
 J\asymp \frac{Y}{Rq_f},
 \qquad J\ge w^{a_\star/2+\delta_0},
\end{equation}
where, for example, one can take \(\delta_0=a_\star/8\) for large \(z\).
The symmetric difference has length \(O(\xi J+1)\).
An enclosing progression piece of length comparable to \(\xi J\) has,
by \cref{lem:small-sieve}, at most \(O(\xi JV_0)\) small-prime survivors.
For fixed \(\xi>0\), its length eventually exceeds \(w^{a_\star/2}\);
this is why the slack in \eqref{eq:variance-common-interval} is retained.
The constant in this survivor bound is independent of sufficiently small
fixed \(\xi\), although the eventual threshold may depend on it.
The reference count changes by \(O(\xi JV_0)\) as well.
Consequently, when \(\xi\) is small relative to a fixed \(\eps>0\), it
suffices to control
\begin{equation}\label{eq:variance-bad-common}
 \left|S-JV_0\right|>\frac{\eps}{2}JV_0,
\end{equation}
where \(S\) counts the small-prime survivors on \(\mathcal J\).

For a prime \(t\le w\), the avoidance condition on \(\mathcal J\) is
\begin{equation}\label{eq:variance-avoidance}
 \begin{cases}
 pq_{f,A}(m_0+Hj)\not\equiv Da_t-A\pmod t,&t\nmid H,\\
 k+pj\not\equiv(a_t-b)\overline{q_f}\pmod t,&t\mid H.
 \end{cases}
\end{equation}
In the first line \(t\nmid D\), so multiplication by \(D\) is reversible.
In the second line \(q_f\) is invertible modulo \(t\).
All factors of \(q_f\) exceed \(w\); hence the primes \(t\le w\) dividing
\(H\) are precisely the small primes dividing \(D\).
Put
\begin{equation}\label{eq:variance-small-products}
 T_0=\prod_{\substack{t\le w\\t\mid H}}t,
 \qquad
 T_1=\prod_{\substack{t\le w\\t\nmid H}}t.
\end{equation}
Thus \((T_1,HT_0)=1\), every prime of \(T_0\) divides \(H\), and
\begin{equation}\label{eq:variance-pattern-size}
 T_0T_1=\exp(O(w))=R^{o(1)}.
\end{equation}

\subsection{Variance of independent classes}

The next second-moment argument is a truncated, two-point form of
mean-one singular-series averaging; compare Gallagher
\cite[equation~(3) and Section~2]{Gallagher1976}.
We include the proof to establish the uniform range of interval lengths
and the stability under changes of measure needed below.

\begin{lemma}\label{lem:random-variance}
Choose, independently at each prime \(t\le w\), one uniformly distributed
forbidden residue class. Let \(S\) count the integers in an interval
\(\mathcal J\) of cardinality \(J\) that avoid all these classes.
For every fixed \(\delta>0\), uniformly for \(J\ge w^\delta\),
\[
 \E S=JV_0,\qquad
 \E(S-JV_0)^2=o((JV_0)^2)
 \quad(w\longrightarrow\infty).
 \]
The assertion has the following stable version. Let \(\mu\) be a finite
positive measure carrying indicators \(I_j\), \(j\in\mathcal J\), and let
\(S=\sum_j I_j\). Suppose, with an absolute fixed implied constant,
\[
 \mu(\Omega)=1+O(\zeta),\qquad
 \int I_j\,d\mu=V_0(1+O(\zeta)),
 \]
and, for distinct \(i,j\), the joint probability has the upper bound
\[
 \int I_iI_j\,d\mu
 \le(1+O(\zeta))\,\Pbb_{\rm ind}(I_i=I_j=1).
 \]
Here \(\Pbb_{\rm ind}\) denotes the independent forbidden-class model
in the first part of the lemma. A zero probability on the right requires
exact zero on the left.
For \(0\le\zeta\le1/2\),
\[
 \int(S-JV_0)^2\,d\mu
 \ll\bigl(\zeta+o(1)\bigr)(JV_0)^2,
 \]
uniformly in the same range of \(J\).
\end{lemma}

\begin{proof}
The single-point survival probability is \(V_0\).
For a nonzero spacing \(h\), the pair probability divided by \(V_0^2\)
equals
\begin{equation}\label{eq:variance-pair-factor}
 2\1_{2\mid h}
 \prod_{2<t\le w}\left(1-\frac1{(t-1)^2}\right)
 \prod_{\substack{2<t\le w\\t\mid h}}
       \left(1+\frac1{t-2}\right).
\end{equation}
At \(t\mid h\) the two positions have one forbidden residue between
them; at \(t\nmid h\) they have two. This proves the formula, including
the parity factor.

Expand the last product in \eqref{eq:variance-pair-factor} over odd
squarefree \(l\), with coefficients
\[
 \lambda_l=\prod_{t\mid l}\frac1{t-2}.
\]
For each fixed \(l\), the triangular average of \(2\1_{2l\mid h}\) is
\[
 \sum_{1\le h<J}\frac{2(J-h)}{J^2}\,2\1_{2l\mid h}
 \longrightarrow\frac1l.
\]
The same average is \(O(1/l)\) for every \(l,J\). Moreover,
\[
 \sum_{\substack{l\ge1\\l\ {\rm odd\ and\ squarefree}}}\frac{\lambda_l}{l}
 =\prod_{t>2}\left(1+\frac1{t(t-2)}\right)<\infty.
\]
Dominated convergence applies even though the available primes increase
with \(w\). The limiting triangular average of
\eqref{eq:variance-pair-factor} is 1, because
\[
 \prod_{t>2}\left(1-\frac1{(t-1)^2}\right)
 \prod_{t>2}\left(1+\frac1{t(t-2)}\right)=1.
\]
The convergence is uniform for \(J\ge w^\delta\): first truncate the
convergent divisor series at fixed \(l\), and then use \(J\to\infty\)
uniformly in this range. The diagonal contribution to
\(\E S^2/(JV_0)^2\) is \(1/(JV_0)=o(1)\), by Mertens' theorem.
This proves the independent-model assertion.

For the stable version, sum the assumed pair upper bounds and the
single-point bounds to obtain
\[
 \int S^2\,d\mu\le
 (1+O(\zeta)+o(1))(JV_0)^2.
\]
Use also the lower bound for \(\int S\,d\mu\) and the upper bound for
\(\mu(\Omega)\) in the identity
\[
 \int(S-JV_0)^2\,d\mu
 =\int S^2\,d\mu-2JV_0\int S\,d\mu+(JV_0)^2\mu(\Omega).
\]
This gives the stated result without normalizing \(\mu\).
\end{proof}

\subsection{The effective-size alternative}

\begin{proposition}\label{prop:structured-variance}
Fix the compact-box parameters, the inverse-estimate parameters and a
relative-error threshold \(\eps>0\). For every fixed \(\sigma_2>0\) there
is a fixed sufficiently large \(C_s\), independent of sufficiently small
fixed \(\xi\), with the following property.
For every candidate box, every choice of its factors other than the
isolated prime, and every rational on its fixed candidate list for which
\eqref{eq:variance-structure-size} holds, suppose
\[
 T>w^{C_s}.
\]
For sufficiently large \(z\), fewer than \(\sigma_2 n_p\) primes
\(p\in\mathcal P\) both align to \(A/D\) and satisfy
\[
 |N_{pq_f}/\mu_{pq_f}-1|>\eps.
\]
The estimate is uniform in the box, the fixed prime choices, the rational
and the forbidden classes. The threshold on \(z\) may depend on fixed
\(\xi\). The fraction is measured against all primes of \(\mathcal P\).
\end{proposition}

\begin{proof}
It is enough to impose the aligned hit at every \(p\) hypothetically and
bound \eqref{eq:variance-bad-common}. The primes dividing \(H\) may be
discarded: there are \(O(1)\) of them in \(\mathcal P\), since
\(H\le R\mathcal Z^{11}\), and their fraction tends to zero.
For \(H\le R^{4/5}\), Brun--Titchmarsh controls each residue pattern
after we fix \(m_0\): the possible primes lie in a progression of
modulus \(HT_0T_1\). The condition that \(T\) is large then gives a
long enough interval in \(m_0\) for its small-prime conditions to have
the required first and second moments. For larger \(H\), we instead
count integer pairs \((p,m_0)\) on a modular hyperbola and then sieve
the \(p\)-coordinate. This gives domination by the uniform distribution
on all small-prime patterns. In both ranges the target is a small
exceptional fraction among all primes of the isolated bin.

\paragraph{The range \(H\le R^{4/5}\).}
As \(p\) varies in \(\mathcal P\), the coordinate \(m_0\) lies in an
enclosing interval \(\mathcal I\) of length
\begin{equation}\label{eq:variance-small-rectangle}
 Y_1\asymp H+\xi|C_0|/R,\qquad (m_0,H)=1.
\end{equation}
For a fixed \(m_0\), possible \(p\)'s lie in an interval of length at most
a constant times
\[
 X_1=\min\left(\xi R,\frac{R^2H}{|C_0|}\right),
\]
using the first entry if \(C_0=0\).
Indeed, \(C_0/p\) must lie in an interval of length \(H\), while its
derivative has magnitude comparable to \(|C_0|/R^2\) on the bin.
This argument is unchanged when \(C_0\) is negative.
Writing \(c=|C_0|/R\), the two cases \(c\le H/\xi\) and \(c>H/\xi\) give
\begin{equation}\label{eq:variance-small-area}
 X_1Y_1\ll\xi RH,\qquad
 Y_1\gg\xi T,\qquad
 X_1\ge R/\mathcal Z^{O(1)}.
\end{equation}
The last inequality follows from \eqref{eq:variance-effective-ratio};
fixed factors depending on \(\xi\) are harmless here.

Specify a unit pattern for \(p\bmod T_1\) and an arbitrary pattern for
\(k\bmod T_0\). Together with fixed \(m_0\), these conditions specify one
unit class for \(p\) modulo \(HT_0T_1\): the condition at \(HT_0\) is
\[
 pm_0\equiv C_0+Hk\pmod{HT_0}.
\]
Both \(m_0\) and the right side are units modulo \(HT_0\), and \(T_1\) is
coprime to this modulus. In particular,
\[
 \varphi(HT_0T_1)=\varphi(H)T_0\varphi(T_1).
\]
Since \(HT_0T_1\le R^{4/5+o(1)}\) and \(X_1=R^{1-o(1)}\),
Brun--Titchmarsh, in the form of \cref{lem:prime-inputs}, bounds the
number of primes for each such choice by
\begin{equation}\label{eq:variance-bt}
 \ll\frac{X_1}{\log R\,\varphi(H)T_0\varphi(T_1)}.
\end{equation}
We may enclose a shorter actual interval by one of the displayed
upper-bound length before applying the theorem.

Define a positive measure on choices consisting of an integer
\(m\in\mathcal I\), \((m,H)=1\), a unit \(p\)-pattern modulo \(T_1\),
and an arbitrary \(k\)-pattern modulo \(T_0\), as follows:
\begin{equation}\label{eq:variance-small-measure}
 \text{weight of each }m=\frac{H}{\varphi(H)Y_1},
 \quad
 p\bmod T_1\text{ and }k\bmod T_0
 \text{ independent and uniform}.
\end{equation}
The number of the latter patterns is \(T_0\varphi(T_1)\).
Dividing \eqref{eq:variance-bt} by
\(n_p\gg\xi R/\log R\), and using
\eqref{eq:variance-small-area}, shows that the fraction of actual primes
whose associated choices lie in any event is at most a fixed absolute
constant times the measure of that event. This constant is independent
of sufficiently small fixed \(\xi\). Each actual prime has one associated
\(m_0\), so this upper count introduces no missing multiplicity.

We verify the hypotheses of the stable part of
\cref{lem:random-variance} for the measure
\eqref{eq:variance-small-measure}.
Take one or two fixed \(j\)'s and let \(\nu_t\in\{1,2\}\) be the number of
their distinct residues modulo \(t\).
At \(t\mid T_0\), reducing the congruence for \(p\) modulo \(t\) shows
that \(p\bmod t\) is already determined by \(m\), independently of the
lift \(k\bmod t\). Uniform \(k\bmod t\) in the second line of
\eqref{eq:variance-avoidance} therefore gives exactly the joint factor
\(1-\nu_t/t\).

At primes \(t\mid T_1\), condition on the unit pattern \(p\bmod T_1\).
The first line of \eqref{eq:variance-avoidance}, as a condition on \(m\),
excludes exactly \(\nu_t\) residue classes, since \(H\) is invertible
there. For a squarefree product \(l\) of these primes, counting the
intersection of its bad classes together with \((m,H)=1\) gives
\begin{equation}\label{eq:variance-small-intersections}
 Y_1\frac{\varphi(H)}H\prod_{t\mid l}\frac{\nu_t}{t}
 +O\left(\tau(H)\prod_{t\mid l}\nu_t\right).
\end{equation}
To see this, expand the coprimality condition over squarefree divisors of
\(H\). Each chosen combination of classes modulo \(l\) combines with
each divisor condition by CRT, because \((l,H)=1\); the interval count
has error \(O(1)\). This gives both the main term and the error in
\eqref{eq:variance-small-intersections}.

Fix a small \(\zeta>0\). Apply \cref{lem:fundamental} at level
\(w^{A_1}\), first choosing the fixed \(A_1\) large enough that its
fundamental-lemma error is smaller than a prescribed constant times
\(\zeta\). If a local product is zero, the conditions are exactly empty
and need no sieve estimate. Otherwise the possible prime 2 has density
at most \(1/2\), and all other densities are at most \(2/t\), as required.
For some fixed \(M>0\), chosen after \(A_1\), the remainder sum is
\(O(\tau(H)w^M)\). Since \(H\ll Y_1\), the elementary divisor and
totient bounds give
\[
 \tau(H)\frac{H}{\varphi(H)}\ll Y_1^{1/4}.
\]
Every nonzero surviving product is \(\gg(\log w)^{-2}\). Thus the
relative remainder is bounded by a constant times
\[
 \begin{aligned}
 \frac{\tau(H)H}{Y_1\varphi(H)}w^M(\log w)^2
 &\ll Y_1^{-3/4}w^M(\log w)^2\\
 &\ll_{\xi}w^{M-3C_s/4}(\log w)^2=o(1),
 \end{aligned}
\]
because \(Y_1\gg\xi T>\xi w^{C_s}\) and we may choose
\(C_s>2M+10\). Thus \(A_1\), then \(M\), then \(C_s\) are fixed
before the eventual choice of \(\xi\). These estimates are uniform in
the pattern, the positions, the interval location and \(H\); a fixed
small \(\xi\) changes only the threshold on \(z\).

After multiplication by the normalization in
\eqref{eq:variance-small-measure}, the single-point and pair probabilities
thus equal their independent-model values with relative error
\(O(\zeta)+o(1)\), whenever those values are nonzero. Factors at \(T_0\)
are exact. The total mass is \(1+o(1)\), either by the same argument with
no bad classes or by the elementary count
\[
 \#\{m\in\mathcal I:(m,H)=1\}
 =Y_1\varphi(H)/H+O(\tau(H)).
\]
The stable variance lemma and Chebyshev's inequality show that the
measure of \eqref{eq:variance-bad-common} is
\(O_\eps(\zeta)+o(1)\). Choose \(\zeta\) sufficiently small in terms of
\(\sigma_2,\eps\) and the domination constant. This proves the desired
prime fraction bound in the first range.

\paragraph{The range \(R^{4/5}<H\le R\mathcal Z^{11}\).}
Partition \(\mathcal P\)'s real bin into equal-length blocks so that
the variation of \(C_0/p\) in each block is at most \(H\).
By \eqref{eq:variance-effective-ratio}, this requires at most
\(\mathcal Z^{O(1)}=R^{o(1)}\) blocks.
For a block of length \(X\), the possible \(m_0\)'s lie in an interval
of length \(Y_2\le3H\). We upper count all pairs in this rectangle with
a full small-prime pattern
\begin{equation}\label{eq:variance-full-pattern}
 p\equiv p_0\pmod{T_0},\quad k\equiv k_0\pmod{T_0},
 \qquad
 p\equiv p_1\pmod{T_1},\quad m_0\equiv m_1\pmod{T_1},
\end{equation}
where \(p_0,p_1\) are units and \(k_0,m_1\) are arbitrary.
The number of such patterns is
\begin{equation}\label{eq:variance-pattern-count}
 M_{\rm pat}=T_0\varphi(T_0)T_1\varphi(T_1)=R^{o(1)}.
\end{equation}
Set \(N=HT_0\). The hyperbola congruence on this rectangle is
\[
 pm_0\equiv C_0+Hk_0\pmod N.
\]
Its right side is a unit by \eqref{eq:variance-unit}.

For any squarefree \(l\le R^{1/10}\) coprime to \(HT_1\), the completion
lemma proved in the appendix, \cref{lem:hyperbola}, gives the count with
\(l\mid p\) as
\begin{equation}\label{eq:variance-hyperbola-count}
 \frac{XY_2\varphi(N)}
 {lT_1^2N^2\varphi(T_0)}
 +O(R^{7/10+o(1)}).
\end{equation}
This is uniform in the block, its location, the pattern and \(l\).
For clarity about the exponent, the lemma's explicit error is bounded by
\[
 N^{1/2+o(1)}T_0^{1/2}
 \left(1+\frac XN+\frac{Y_2}N\right).
\]
Since \(R^{4/5}<N\le R^{1+o(1)}\), \(X\ll R\),
\(Y_2\le3H\le3N\), and \(T_0=R^{o(1)}\), this is in fact
\(O(R^{3/5+o(1)})\). We use the weaker error in
\eqref{eq:variance-hyperbola-count} to leave slack in subsequent sums.

Sieve the \(p\)-coordinate of this nonnegative sequence of pairs for an
upper bound on primes. Exclude primes at most \(R^{c_6}\) not dividing
\(HT_1\), where fixed \(c_6>0\) is sufficiently small, and use level
\(R^{1/10}\) in \cref{lem:fundamental}.
The intersection density supplied by
\eqref{eq:variance-hyperbola-count} is \(1/l\).
The primes dividing \(HT_1\) are already excluded from the \(p\)-coordinate
by the unit conditions. The main sieve product is at most
\begin{equation}\label{eq:variance-large-sieve-product}
 \ll\frac1{\log R}\frac H{\varphi(H)}\frac{T_1}{\varphi(T_1)}.
\end{equation}
Including factors of \(H\) larger than the sieve cutoff only increases
this upper bound.
The fundamental-lemma parameter is \((1/10)/c_6\), so a fixed small
choice of \(c_6\) suffices.

The weighted sum of the errors in
\eqref{eq:variance-hyperbola-count} over \(l\le R^{1/10}\) is
\(R^{4/5+o(1)}\), and in particular is \(O(R^{9/10})\) for large \(z\).
Because every prime of \(T_0\) divides \(H\),
\(\varphi(N)=\varphi(H)T_0\). Multiplying the main term at \(l=1\) by
\eqref{eq:variance-large-sieve-product} and using \(Y_2\le3H\) therefore
gives, for each pattern and block,
\begin{equation}\label{eq:variance-pattern-primes}
 \ll \frac{X}{\log R\,T_0\varphi(T_0)T_1\varphi(T_1)}
 +O(R^{9/10}).
\end{equation}
There are \(R^{o(1)}\) blocks and \(R^{o(1)}\) patterns, so all error
terms together are
\[
 R^{9/10+o(1)}=o(n_p).
\]
Here division by \(n_p\asymp\xi R/\log R\) is harmless for each fixed
\(\xi>0\). Summing the main term over the blocks uses
\(\sum X\asymp\xi R\). Thus the fraction of actual primes associated to
any set of bad full patterns is bounded by a fixed constant times its
fraction among all \(M_{\rm pat}\) patterns, plus \(o(1)\).
The bound counts pairs, which is sufficient: every actual prime has
exactly one \(m_0\) from \eqref{eq:variance-m-coordinate}.

Under independent uniform choices in \eqref{eq:variance-full-pattern},
the forbidden \(j\)-classes in \eqref{eq:variance-avoidance} are exactly
independent and uniform. Conditional on the unit patterns, use uniform
\(k_0\bmod t\) at each \(t\mid T_0\); at the other primes use uniform
\(m_1\bmod t\). The coefficients of these variables are units.
CRT supplies independence across the primes.
The event \eqref{eq:variance-bad-common} depends only on this full
pattern and on the already fixed \(\mathcal J,q_f,A,D,b\).
By \cref{lem:random-variance}, its uniform-pattern fraction is \(o(1)\).
The domination just proved gives the required prime fraction bound in
the second range.

Finally, the first range chose \(\zeta\), then \(A_1\), then \(C_s\);
the second range needs no further increase of \(C_s\).
These choices are independent of sufficiently small fixed \(\xi\).
Choose \(\xi\) small enough for the earlier interval freezing and then
let \(z\) tend to infinity. This proves all asserted uniformities.
\end{proof}

\subsection{Hard endpoints and global eligibility}

\begin{corollary}\label{cor:hard-endpoints}
Fix the compact range and \(\eps>0\).
For every prescribed positive \(\rho\), the inverse exception parameter
and then the variance exception parameter can be chosen so that, for a
fixed \(C_s\) independent of sufficiently small fixed \(\xi\), the
following holds for all sufficiently large \(z\).
Consider the actual included regular compact endpoints covered by the
full boxes of \cref{prop:regular-boxes}. Except for a set of these
endpoints of total harmonic mass at most \(\rho B^{-2}\), every such
bad endpoint
\[
 |N_d/\mu_d-1|>\eps
\]
has a witnessing candidate and a rational on that candidate's fixed
list for which the parent is structured, the isolated prime aligns, and
\begin{equation}\label{eq:variance-hard}
 Dq_{f,E}\le w^{C_s},
 \qquad |C_0|/R\le w^{C_s}.
\end{equation}
For each such endpoint all its prime factors with exponent
\(x(t)>C_s\) align to this same rational, including the isolated prime.
The rational satisfies the global bounds
\begin{equation}\label{eq:variance-global-height}
 D\le w^{C_s},
 \qquad |A|\le Yw^{C_s+2}.
\end{equation}
The enlarged full boxes are used only to dominate measure: the assertion
concerns the actual included endpoints, rather than every tuple in
those enlarged boxes.
The exceptional mass can be bounded before choosing \(\xi\); only the
eventual threshold on \(z\) is allowed to depend on fixed \(\xi\).
\end{corollary}

\begin{proof}
Every actual included bad endpoint under consideration has a witness
from the bounded candidate set, by
\cref{lem:edge-witness,prop:regular-boxes}.
For each candidate, \cref{prop:inverse} costs at most a prescribed
\(\sigma\) fraction of its full harmonic box and supplies a fixed list
of at most \(F\) rationals. Both \(F\) and the list definition are
independent of the choice of the isolated prime.

To apply \cref{prop:structured-variance}, fix all the other prime choices
and one rational on this list. Then \(H,C_0,T\) are fixed as \(p\) varies.
If the structured-parent condition is possible in this slice,
\eqref{eq:variance-structure-size} holds throughout the fixed slice.
When \(T>w^{C_s}\), the proposition bounds all bad aligning \(p\)'s,
and hence also those that supply the particular witness.
Since the isolated bin has multiplicity one, its prime choice is
independent of all the other bin choices in full-box measure.
The weights \(1/p\) vary by a bounded factor within this bin, uniformly
for sufficiently small \(\xi\). Therefore the count bound
\(\sigma_2n_p\) converts to at most \(C\sigma_2\) of its harmonic
measure, with fixed \(C\).

Union over the list and over the at most \(H_e\) candidates bounds the
discarded fraction of each full box by
\[
 H_e(\sigma+CF\sigma_2).
\]
The total enlarged-box mass is at most \(C'_K B^{-2}\), where \(C'_K\)
is independent of sufficiently small \(\xi\), by
\cref{prop:regular-boxes}. First choose \(\sigma\) sufficiently small;
this fixes \(F\). Then choose \(\sigma_2\) sufficiently small in terms
of \(F\) and \(\rho\), and finally choose \(C_s\) from
\cref{prop:structured-variance}. This makes the discarded mass at most
\(\rho B^{-2}\), without multiplying by the number of boxes.
For every remaining bad endpoint, some structured rational has
\(T\le w^{C_s}\), which is \eqref{eq:variance-hard}.

The bound on \(H\) controls the nonaligned factors. Each nonaligned
prime factor of \(q_f\) divides \(q_{f,E}\le w^{C_s}\), so it is at
most \(w^{C_s}\).
The isolated prime already aligns by the structured-parent conclusion;
its exponent tends to infinity, and in particular exceeds \(C_s\).
This proves the alignment assertion.
The bound on \(|C_0|/R\) controls the numerator of the rational on a
scale common to all boxes. We have \(b\le q_f\), \(q_{f,A}\le q_f\),
\(R\le p\), and
\(pq_f=d\le Y\), since these are included endpoints with \(J_d>1\).
Consequently,
\[
 |A|\le Dq_f+|C_0|q_{f,A}
 \le w^{C_s}q_f(1+R)
 \le 2Yw^{C_s}.
\]
This implies the second inequality of
\eqref{eq:variance-global-height} for large \(z\); the first follows
directly from \eqref{eq:variance-hard}.
\end{proof}
